% §14.7 Bibliographic Notes
%
% Monograph register: primary source (Towell BHF manuscript),
% universal-hashing foundation (Carter-Wegman), modern perfect hash
% (BDZ), comparison structures (cuckoo, quotient). Light compression to
% ch13 density. Cites towell-bernoulli-hash-function, carter-wegman-1979,
% belazzougui-2009, fan-cuckoo-2014, bender-quotient-2012.
% DEFINED: sec:bib-notes-ch14.

\section*{Bibliographic Notes}
\addcontentsline{toc}{section}{Bibliographic Notes}
\label{sec:bib-notes-ch14}

\paragraph{The BHF construction.}
The construction of \Cref{def:bhf} and the optimality result of
\Cref{thm:bhf-optimality} follow
\textcite{towell-bernoulli-hash-function}, ``The Bernoulli Hash
Function: Optimal Bernoulli Sets and Maps,'' which introduced the BHF as
the constructive complement to the chapter-8 floor. The manuscript works
in greater generality, extending the membership construction of
\Cref{sec:bhf-construction} to approximate maps (each slot storing a
key-value pair rather than a fingerprint); the space and FPR analyses here
specialize that framework to set membership. A reader wanting the algebraic
generalization, with the BHF as an instance of the optimal Bernoulli map,
should treat the manuscript as the next reading. The perfect hash filter of
chapter~15 is a distinct, complementary construction of the same set and map
ADT, not an instance of the BHF: it builds by off-the-shelf perfect hashing
at a constant $\log_2 e$ above the floor the BHF reaches by salt search.

\paragraph{The universal-hashing foundation.}
The fingerprint hash $g$ of \Cref{prop:bhf-fpr} inherits its uniformity
from the universal hashing of \textcite{carter-wegman-1979}, ``Universal
Classes of Hash Functions,'' the first family with provable pairwise
independence (now $2$-universal), bounding collision probabilities to
$1/|\text{codomain}|$ over the family's random selection. The BHF needs
only pairwise independence, so the later $k$-universal refinements
(Wegman-Carter 1981, Siegel 1989) are not required, and Carter-Wegman
remains the standard reference for the result the analysis depends on.

\paragraph{The modern perfect hash.}
The sublinear bound of \Cref{sec:space-analysis} and the
linear-expected-time algorithms of \Cref{sec:construction-time} trace to
a long line culminating in \textcite{belazzougui-2009}, ``Hash,
Displace, and Compress'' (the BDZ scheme), which achieves $O(n)$
expected construction with output in $O(n \log \log \log n / \log n)$ at
under $3$ bits per element for $n$ in the millions, combining
hash-and-displace placement with entropy-efficient compression of the
displacements. The chapter quotes its asymptotic bound as a black box.
The antecedents BDZ improves on are Fredman-Komlos-Szemeredi (1984, the
classical $O(n \log n)$-bit bound), Hagerup-Tholey (2001, the first
$O(n)$-bit construction), and the MWHC and BBHash displacement families.

\paragraph{Comparison structures.}
Two contemporary designs occupy spaces adjacent to the BHF.
\textcite{fan-cuckoo-2014}'s cuckoo filter stores fingerprints in a
cuckoo-hashed, cache-friendly layout, reaching about $1.05$ to $1.3$
times the floor (against the Bloom filter's $1.44$) while supporting
amortized $O(1)$ insertion and deletion; it recovers about half the
Bloom overhead while remaining streamable.
\textcite{bender-quotient-2012}'s quotient filter stores fingerprints in
a quotient-style addressing scheme with similar trade-offs and streaming
properties. Neither matches the BHF's asymptotic optimality, but both
narrow the Bloom gap and extend the streaming case beyond what the BHF
supports, making them the practical middle ground when a deployment
needs both space efficiency and streaming updates.
